%% LyX 2.3.4.2 created this file.  For more info, see http://www.lyx.org/.
%% Do not edit unless you really know what you are doing.
\documentclass[english]{article}
\usepackage[T1]{fontenc}
\usepackage[utf8]{inputenc}
\usepackage{amsmath}
\usepackage{graphicx}

% (refstyle removed — amsmath provides \eqref; \subsecref usages replaced with vanilla \ref)
\providecommand\subsecref[1]{Section~\ref{subsec:#1}}
\providecommand\secref[1]{Section~\ref{#1}}
\providecommand\figref[1]{Figure~\ref{#1}}
% Render \cite as [TBD: ...] without a bibliography
\renewcommand\cite[1]{[\textit{cite: #1}]}

\usepackage{babel}
\begin{document}
%\title{Electro-chemo-mechanics of Dendrite Growth in Li-Metal Solid-state %Battery}

\author{Howard Tu}

\title{A Coupled Electro-Chemo-Mechanical Framework for Lithium-Metal Solid-State Batteries: Growth, Plasticity, Diffusion, and Contact}

\maketitle

\section{Introduction}
\label{sec:introduction}

Solid-state lithium-metal batteries (SSBs) promise significant gains in specific energy and operational safety relative to the liquid-electrolyte status quo. Replacing the flammable organic electrolyte with a ceramic or glassy solid electrolyte (SE) eliminates the leakage and flammability hazards associated with conventional Li-ion cells, while enabling the use of a high-capacity Li-metal anode that has been precluded in liquid systems by dendrite-driven failure~\cite{TBD:JanekZeier,TBD:Famprikis}. Practical deployment of SSBs is, however, gated by mechanical and electrochemical failure mechanisms whose origin is fundamentally a coupling between transport, deformation, and contact: contact-pressure build-up at the SE/Li interface during plating can fracture the SE; dendritic Li penetration along grain boundaries shorts the cell; voiding under stripping disconnects the anode; and stack-pressure variations modulate all three. None of these failure modes is captured by stand-alone electrochemical or stand-alone mechanical models — they require a unified treatment.

The mechanical behavior of Li metal is the central piece of physics tying these failure modes together. Even at room temperature, polycrystalline Li exhibits viscoplastic creep at stress levels and time scales relevant to galvanostatic cycling~\cite{TBD:LeSar,TBD:Masias,TBD:LePage}. Plastic flow simultaneously relaxes contact pressure (suppressing some failure modes) and supplies the bulk source of metal needed to feed dendrite growth (driving others). Because the kinematics of growth at a Li-metal/SE interface are not those of bulk plasticity --- new material appears at a moving interface, and that material must be elastically and plastically accommodated by the surroundings --- a simulation framework for SSBs must combine: (a) a multiplicative kinematic decomposition of the deformation gradient that distinguishes elastic, plastic, and growth contributions; (b) a continuum representation of $\mathrm{Li}^+$ transport in the SE coupled to neutral Li transport in the interphase; (c) a contact-mechanics treatment of the SE/anode interface that handles non-matching meshes, mortar-FEM constraints, and pressure-modulated kinetics; and (d) a free-energy / chemical-potential formulation that couples (a)--(c) thermodynamically.

Several research groups have addressed subsets of this problem. Bucci et al.~\cite{TBD:Bucci} and Bistri et al.~\cite{TBD:Bistri} developed pioneering coupled mechanical-electrochemical models for SSBs and reproduced contact-pressure-driven failure scenarios. McMeeking and co-workers~\cite{TBD:McMeeking,TBD:Purkayastha} built phase-field models of dendrite extension that capture the role of stack pressure on filament morphology. The Anand group's viscoplastic Li constitutive model~\cite{TBD:LiAnand} provides the closure relation for bulk Li flow under cell-relevant conditions, and works including~\cite{TBD:Bistri2024,TBD:Bonnin} have applied this constitutive model to specific cell geometries. These efforts have advanced understanding of individual failure modes but have not been unified within a single framework that exposes all of (a)--(d) as user-controllable physics modules within a common finite-element implementation.

This paper presents such a unified framework, implemented as a custom application within the MOOSE finite-element platform. The framework couples a multiplicative-decomposition kinematics for elastic-plastic-growth deformation (\secref{sec:methodology}, in particular \subsecref{Free-growth}) with anisotropic Li diffusion in the interphase, Nernst-Planck $\mathrm{Li}^+$ migration in the SE (\secref{subsec:li_migration_se}), mortar-FEM contact (\secref{subsec:contact}), and a thermodynamically-consistent chemical-potential closure (\secref{subsec:free_energies}). The numerical implementation is described in \secref{subsec:numerical_implementation}. We exercise the framework through a hierarchy of benchmarks (\secref{sec:benchmarks}) that build up from the constitutive Li-metal creep response, through free and confined growth in the interphase, to fully-coupled $\mathrm{Li}^+$/$\mathrm{Li}$ transport. We then deploy it on five plating / stripping scenarios (\secref{sec:plating_stripping}), including validation against McMeeking-type cell-level experiments (\subsecref{mcmeeking}) and dendrite-onset phenomenology (\subsecref{dendrite}). Conclusions and avenues for extension are summarized in \secref{sec:conclusions}.



\section{Methodology}
\label{sec:methodology}

\subsection{Kinematics of Deformation}
Consider a macroscopically-homogeneous body \textbf{B} with the region
of space it occupies in a fixed reference configuration, and denote
by \textbf{X} an arbitrary material point of \textbf{B}. A motion
of \textbf{B} is then a smooth one-to-one mapping $\mathbf{x}=\mathbf{\mathbf{\chi}}(\mathbf{X},t)$
with\emph{ deformation gradient, velocity, and velocity gradient}
given by
\begin{equation}
\mathbf{F}=\nabla\mathbf{\chi},\mathbf{v}=\mathbf{\dot{{\mathbf{\chi}}}},\mathbf{L}=\mathrm{grad}\mathbf{v}=\mathbf{\dot{F}F^{-1}}\label{eq:def_grad_defn}
\end{equation}

Assume that $J=\mathrm{det}\mathbf{F}>0$, where $J=\frac{dV}{dV_{0}}$.
Following modern developments in large-deformation plasticity and
growth we regard the growth of an interphase layer on the multiplicative
decomposition of the deformation gradient, 

\begin{equation}
\mathbf{F=F^{e}F^{p}F^{g}}\label{eq:def_grad_decomp}
\end{equation}

Here
\begin{enumerate}
\item $\mathbf{F^{g}(X},t)$ represents the local distortion of the material
neighborhood of $\mathbf{X}$ due to growth by deposition of Li metal;
\item $\mathbf{F^{p}(X},t)$ represents the constant volume plastic (viscoplastic)
deformation of the deposited Li;
\item $\mathbf{F^{e}(X},t)$ represents the elastic distortion that occurs
due to the subsequent stretching and rotation. 
\end{enumerate}
$J$ now follows the multiplicative decomposition of the deformation
gradient according to 

\begin{equation}
J=J^{e}J^{p}J^{g}\label{eq:J_decomp}
\end{equation}

where $J^{e}=\mathrm{det}\mathbf{F}^{e}>0$, $J^{p}=\mathrm{det}\mathbf{F}^{p}=1$
(constant volume deformation), and $J^{g}=\mathrm{det}\mathbf{F}^{g}>0$
. We further define the elastic, plastic and growth stretch and spin
tensors through 

\begin{equation}
\begin{array}{cc}
\mathbf{D}^{e}=\mathrm{symm}\mathbf{L}^{e}, & \mathbf{W}^{e}=\mathrm{skew}\mathbf{L}^{e}\\
\mathbf{D}^{p}=\mathrm{symm}\mathbf{L}^{p}, & \mathbf{W}^{p}=\mathrm{skew}\mathbf{L}^{p}\\
\mathbf{D}^{g}=\mathrm{symm}\mathbf{L}^{g}, & \mathbf{W}^{g}=\mathrm{skew}\mathbf{L}^{g}
\end{array}\label{eq:stretch_spin_tensors_defn}
\end{equation}

where $\mathbf{L=L}^{e}+\mathbf{F}^{e}\mathbf{L}^{p}\mathbf{F}^{e-1}+(\mathbf{F}^{e}\mathbf{F}^{p})\mathbf{L}^{g}(\mathbf{F}^{e}\mathbf{F}^{p})^{-1}$ and
$\mathbf{L}^{e}=\mathbf{\dot{F}}^{e}\mathbf{F}^{e-1},$$\mathbf{L}^{p}=\mathbf{\dot{F}}^{p}\mathbf{F}^{p-1}$,
$\mathbf{L}^{g}=\mathbf{\dot{F}}^{g}\mathbf{F}^{g-1}$. We further assume
that the plastic spin and growth spin tensors are identically zero. 

\subsection{Growth Deformation in Interphase Layer}

We introduce a set of orthonormal unit vectors$\{\mathbf{m}_{Ri}\vert i=1,2,3\}$
in the reference body and consider a particular class of growth tensors
$\mathbf{F}^{g}$ defined by 

\begin{equation}
\mathbf{F}^{g}=\sum_{i=1}^{3}\lambda_{i}^{g}(\mathbf{m}_{Ri}\otimes\mathbf{m}_{Ri})\label{eq:Fg_defn}
\end{equation}

where $\{\lambda_{i}^{g}>0\vert i=1,2,3\}$ are the corresponding
growth stretches which satisfy 

\begin{equation}
\prod_{i=1}^{3}\lambda_{i}^{g}=J^{g}\label{eq:Jg_defn}
\end{equation}

Since the growth is assumed to have no spin we can define from \ref{eq:Fg_defn} 

\begin{equation}
\begin{aligned}\mathbf{\dot{F}}^{g}=\sum_{i=1}^{3}\dot{\lambda}_{i}^{g}(\mathbf{m}_{Ri}\otimes\mathbf{m}_{Ri}), \: \mathrm{and} \: & \mathbf{F}^{g-1}=\sum_{i=1}^{3}\lambda_{i}^{g-1}(\mathbf{m}_{Ri}\otimes\mathbf{m}_{Ri})\end{aligned}
\label{eq:growth_def_grad_rate}
\end{equation}

From \ref{eq:growth_def_grad_rate} we can now define the growth stretch
rate as 

\begin{equation}
\begin{split}\begin{split}\mathbf{D}^{g}=\sum_{i=1}^{3}\dot{\lambda}_{i}^{g}\lambda_{i}^{g-1} & (\mathbf{m}_{Ri}\otimes\mathbf{m}_{Ri})\end{split},
\: \mathrm{and} \: & \mathbf{W}^{g}=0\end{split}
\label{eq:growth_stretch_rate_defn}
\end{equation}

Define $\dot{\epsilon}_{i}^{g}=\dot{\lambda}_{i}^{g}\lambda_{i}^{g-1}$, and $\dot{\epsilon}^{g}$ the volumetric growth strain rate: $\dot{\epsilon}^{g} \equiv \mathrm{tr}\mathbf{D}^{g}$. Therefore: $\dot{\epsilon}^{g}=\mathrm{tr}\mathbf{D}^{g}=\sum_{i=1}^{3}\dot{\epsilon}_{i}^{g}=\dot{J}^{g}J^{g-1}$.
The volumetric growth strain $\epsilon^{g}$ is obtained by integrating $\dot{\epsilon}^{g}$ over time:

\begin{equation}
\epsilon^{g}=\ln J^{g}\label{eq:vol_growth_strain}
\end{equation}

If $\dot{\epsilon}^{g}>0$ material growth or Li is deposited and
$\dot{\epsilon}^{g}<0$ materials shrinks or Li is stripped. 

To allow for different growth rates along the preferred directions $\{\mathbf{m}_{Ri}\vert i=1,2,3\}$,
we introduce a triplet of growth-weight parameters $\{\alpha_{i}\in[0,1]\vert i=1,2,3\}$
and assume that the growth strain rate in each direction is a fraction
of the volumetric growth strain rate 

\begin{equation}
\dot{\epsilon}_{i}^{g}=\alpha_{i}\dot{\epsilon}^{g}\label{eq:growth_weights_defn}
\end{equation}
The growth-weight parameters $\alpha_{i}$ must satisfy the constraint 

\begin{equation}
\sum_{i=1}^{3}\alpha_{i}=1\label{eq:growth_weight_constraint}
\end{equation}

Now using the above equations we note that 

\begin{equation}
\dot{\lambda}_{i}^{g}\lambda_{i}^{g-1}=\alpha_{i}\dot{J}^{g}J^{g-1}\label{eq:growth_weight_Jg}
\end{equation}

which can now be integrated to give 

\begin{equation}
\lambda_{i}^{g}=\left(J^{g}\right)^{\alpha_{i}}\label{eq:individual_growth_strethc}
\end{equation}

Finally, all the above arguments can now be condensed to the rate
form for the evolution of $\mathbf{F}^{g}$ as 

\begin{equation}
\dot{\mathbf{F}}^{g}=\mathbf{D}^{g}\mathbf{F}^{g}\label{eq:Fg_evolution}
\end{equation}

with 

\begin{equation}
\mathbf{D}^{g}=\dot{\epsilon}^{g}\mathbf{S}^{g}\label{eq:Dg_eqn}
\end{equation}

where $\mathbf{S}^{g}$ is the principal growth tensor and defined
as 

\begin{equation}
\mathbf{S}^{g}=\sum_{i=1}^{3}\alpha_{i}(\mathbf{m}_{Ri}\otimes\mathbf{m}_{Ri})\label{eq:Sg_reference}
\end{equation}


\subsection{Non-fixed reference vector}

We can now extend this analysis to a non-fixed reference vector defined
in the intermediate configuration of $\mathbf{F}^{g}$. We replace
$\mathbf{m}_{Ri}$ the reference vector by $\mathbf{m}_{i}^{g}$ which
are unit vectors in the intermediate space of $\boldsymbol{F}_{g}$. $\mathbf{m}_{i}^{g}$ is defined as 
\begin{equation}
    \mathbf{m}_{i}^{g} = \frac{(\mathbf{F}^e\mathbf{F}^p)^{-1} \mathbf{m}_{i}}{\lVert (\mathbf{F}^e\mathbf{F}^p)^{-1} \mathbf{m}_{i} \rVert} \label{eq:principal_vector_transformation}
\end{equation}
where $\mathbf{m}_{i}$ are orthonormal vectors in the deformed configuration. There is no guarantee that the vectors $\mathbf{m}_{i}^{g}$ are orthonormal. So in order to have a set of order to have a set of vectors orthonormal vectors, we choose $\mathbf{m}_{1}^{g}$ as one of the principal directions and then define two orthogonal vectors $\mathbf{m}_{2}^{g}$ and $\mathbf{m}_{3}^{g}$. Thus the transformation described in \ref{eq:principal_vector_transformation} is only for the chosen principal direction. 


Thus the principal growth tensor in \ref{eq:Sg_reference} is now
defined in the intermediate space as 

\begin{equation}
\mathbf{S}^{g}=\sum_{i=1}^{3}\alpha_{i}(\mathbf{m}_{i}^{g}\otimes\mathbf{m}_{i}^{g})\label{eq:principal_growth_tensor_interme}
\end{equation}

\begin{description}
\item [{Warning:}] With a non-fixed reference vector the direct formulation
for $\mathrm{\boldsymbol{F}}_{g}$ \eqref{Fg_defn} cannot be used
any longer and we have to rely on the incremental formulation. 
\end{description}

\subsubsection*{Axial (Fiber) growth\label{subsec:Axial-(Fiber)-growth}}

This is a specialization of the equation for growth along a particular
axis. Instead of defining all three principal directions, we define
a single principal (fiber) direction $\mathrm{\boldsymbol{m}}_{1}$
along which growth occurs. For this form of growth, 

\[
\begin{aligned}\alpha_{1}=1, & \alpha_{2}=\alpha_{3}=0,\end{aligned}
\mathrm{so\;that}
\]

\begin{equation}
\mathbf{S}^{g}=\left(\mathrm{\boldsymbol{m}}_{1}\otimes\mathrm{\boldsymbol{m}}_{1}\right)\label{eq:axial_Sg_fiber_direction}
\end{equation}


\subsubsection*{Areal growth (Growth perpendicular to a single principal direction)}

A similar specilization to the previous section \subsecref{Axial-(Fiber)-growth}
can be applied to growth perpendicular to a single principal direction
$\mathrm{\boldsymbol{m}}_{1}$. For this form of growth, 

\[
\begin{aligned}\alpha_{1}=0, & \alpha_{2}=\alpha_{3}=\frac{1}{2}\end{aligned}
,\mathrm{so\;that}
\]

\begin{equation}
\mathrm{\boldsymbol{S}}_{g}=\frac{1}{2} (\boldsymbol{1}-\mathrm{\boldsymbol{m}_{1}}\otimes\mathrm{\boldsymbol{m}_{1}})\label{eq:Sg_area_growth}
\end{equation}

In order to solve for the evolution of $\mathrm{\boldsymbol{F}}_{g}$
\eqref{Fg_evolution} we assume that at the start of deformation the
initial condition for $\mathrm{\boldsymbol{F}}_{g}(\mathrm{\boldsymbol{X}},t=0)=\mathrm{\boldsymbol{1}}$.
Therefore, 

\begin{equation}
\mathrm{\boldsymbol{F}_{g}(t+\Delta t})=\exp(\mathrm{\boldsymbol{D}_{g}\Delta t})\mathrm{\boldsymbol{F}_{g}(t)}\label{eq:F_g_actual_equation}
\end{equation}


\subsubsection*{Constitutive equation for swelling }

The swelling volume change $J^{g}$ is related to the Li concentration
in the interphase layer through a linear constitutive equation 

\begin{equation}
J^{g}=1+\Omega(c_{R}-c_{R0})\label{eq:Jg_constitutive}
\end{equation}

where $c_{R}$is the concentration, $c_{R0}$ is a reference concentration,
with $\Omega>0$ representing the \emph{molar volume of Li. }From
\eqref{vol_growth_strain} and \eqref{Jg_constitutive} we now have 

\begin{equation}
\dot{J}^{g}J^{g}=\Theta(c_{R})\dot{c}_{R}\label{eq:Jg_evolution_const}
\end{equation}

where

\begin{equation}
\Theta(c_{R})=\frac{\Omega}{1+\Omega c_{R}}\label{eq:theta_const}
\end{equation}

we obtain the volumetric growth strain rate as 

\begin{equation}
\dot{\epsilon}^{g}=\Theta(c_{R})\dot{c}_{R}\label{eq:vol_sgrowth_train_rate_const}
\end{equation}

From \eqref{vol_sgrowth_train_rate_const} and \eqref{theta_const}
the growth stretch rate of \eqref{Dg_eqn} can now be written as 

\begin{equation}
\mathrm{\boldsymbol{D}}^{g}=\Theta(c_{R})\dot{c}_{R}\mathrm{\boldsymbol{S}_{g}}\label{eq:Dg_const_defn}
\end{equation}


\subsection{Li Diffusion in Interphase Layer}

The Li flux in the reference frame is presumed to obey the constitutive
equation 

\begin{equation}
\mathrm{\boldsymbol{j}_{R}}=-\left[\sum_{i=1}^{3}m_{i}(c_{R})\mathrm{\boldsymbol{m}}_{Ri}\otimes\mathrm{\boldsymbol{m}}_{Ri}\right]\nabla\mu\label{eq:reference_Li_flux}
\end{equation}

where $m_{i}$ is a scalar mobility of Li in the $\mathrm{\boldsymbol{m}}_{R}$-direction.
The spatial flux is related to the reference flux by 

\begin{equation}
\mathrm{\boldsymbol{j}}_{R}=J\mathrm{\boldsymbol{F}}^{-1}\mathrm{\boldsymbol{j}}\label{eq:ref_flux_spatial_flux}
\end{equation}

and 

\begin{equation}
\nabla\mu=\mathrm{\boldsymbol{F}}^{T}\mathrm{grad}\mu\label{eq:gradient_chem_pot_ref_spat}
\end{equation}

From \eqref{reference_Li_flux}, \eqref{ref_flux_spatial_flux} and
\eqref{gradient_chem_pot_ref_spat} we can now write the spatial flux
as 

\begin{equation}
\mathrm{\boldsymbol{j}}=-\left[J^{-1}\sum_{i=1}^{3}m_{i}(c_{R})\mathrm{\boldsymbol{m}}_{i}\otimes\mathrm{\boldsymbol{m}}_{i}\right]\mathrm{grad}\mu\label{eq:spatial_flux}
\end{equation}

with $\mathrm{\boldsymbol{m}}_{i}=\mathrm{\boldsymbol{F}}\mathrm{\boldsymbol{m}}_{Ri}$
and $\mu$ is the chemical potential. 

The chemical potential can take many forms 
\begin{enumerate}
\item $\mu=\mu_{0}+R\vartheta\log\left(\frac{\gamma c_{R}}{c_{R0}}\right)+\Omega(\frac{1}{2}\mathrm{\boldsymbol{E}}^{e}:\mathcal{C}\mathrm{\boldsymbol{E}}^{e})-\omega$,
where $\mathrm{\boldsymbol{E}^{e}}$ is the logarithmic elastic strain,
$\mathcal{C}$ is the isotropic elasticity tensor, $\vartheta$ is
the temperature, $\gamma$ is the activity coefficient, $R$ is the
universal gas constant, and $\omega$ is the growth chemical potential 
\item $\mu=R\vartheta\ln\left(\frac{\bar{c}}{1-\bar{c}}\right)+\Omega(\frac{1}{2}\mathrm{\boldsymbol{E}}^{e}:\mathcal{C}\mathrm{\boldsymbol{E}}^{e})-\omega$
\end{enumerate}
For the model described above 

\begin{equation}
\omega=\Omega\mathrm{\boldsymbol{M}^{g}:\mathrm{\boldsymbol{S}}^{g}}\label{eq:growth_chemical_potential}
\end{equation}

where 

\begin{equation}
\mathrm{\boldsymbol{M}}^{g}=\mathrm{\boldsymbol{F}}^{pT}\mathrm{\boldsymbol{M}}^{e}\mathrm{\boldsymbol{F}}^{p-T}\label{eq:growth_stress_eqn}
\end{equation}

and 

\begin{equation}
\mathrm{\boldsymbol{M}}^{e}=J^{e}\mathrm{\boldsymbol{R}}^{eT}\mathrm{\boldsymbol{\sigma}}\mathrm{\boldsymbol{R}}^{e}\label{eq:Mandel_cauchy_stress_rel}
\end{equation}

is the Mandel stress. The description of these quantities and the
thermodynamic derivation of these quantities is beyond the scope of
this document. 

In Li dominated problems under small stress limits, the growth stress
and the elastic stresses are expected to be small, so both the growth
and the elastic terms can be ignored for the chemical potential. 
Once the elastic energy and the growth energy contributions to the chemical potential are
ignored we can now readily write the flux in standard Fickian form as

\begin{equation}
\mathbf{j} = -\mathbf{D}(\mathbf{x}) \cdot \nabla c \label{eq:aniso_diffusion_flux}
\end{equation}
where $\mathbf{D}(\mathbf{x}) = (R\vartheta/c)\, \mathbf{M}$, with $\mathbf{M}$ being the mobility tensor (first bracketed term in \eqref{eq:reference_Li_flux}).

Combining \eqref{eq:aniso_diffusion_flux} with mass conservation $\partial c/\partial t + \nabla\cdot\mathbf{j} = 0$ yields the anisotropic diffusion equation in divergence form

\begin{equation}
    \frac{\partial c}{\partial t} = \nabla \cdot \left[\mathbf{D}(\mathbf{x})\, \nabla c\right] \label{eq:aniso_diffusion_pde}
\end{equation}

For numerical implementation the right-hand side may be expanded by the product rule as $\nabla\cdot\mathbf{D}\cdot\nabla c + \mathbf{D}{:}\nabla\nabla c$, but the conservative form \eqref{eq:aniso_diffusion_pde} is preferred for finite-element weak forms.

\subsection{$Li^+$ migration in the Solid Electrolyte}
\label{subsec:li_migration_se}

In the solid electrolyte the relevant transported species is $\mathrm{Li}^+$. Charge transport is governed by the Nernst-Planck equation in the limit of dilute solution and small concentration variation, which writes the molar flux of $\mathrm{Li}^+$ as the sum of a diffusive Fickian term and an electromigration term driven by the gradient of the electrostatic potential $\phi$:
\begin{equation}
    \mathbf{j}_{\mathrm{Li}^+} = -D_{\mathrm{Li}^+}\, \nabla c_{\mathrm{Li}^+} - \frac{z_{\mathrm{Li}^+} F}{R\vartheta}\, D_{\mathrm{Li}^+}\, c_{\mathrm{Li}^+}\, \nabla \phi
    \label{eq:nernst_planck_flux}
\end{equation}
where $D_{\mathrm{Li}^+}$ is the $\mathrm{Li}^+$ diffusivity in the SE, $z_{\mathrm{Li}^+} = +1$ is the charge number, $F$ is Faraday's constant, $R$ is the universal gas constant, and $\vartheta$ is the absolute temperature. The Einstein relation $D_{\mathrm{Li}^+} = \mu_{\mathrm{Li}^+} R\vartheta / |z_{\mathrm{Li}^+}|F$ with mobility $\mu_{\mathrm{Li}^+}$ is implicit in this form.

The current density is related to the ionic flux by $\mathbf{i} = z_{\mathrm{Li}^+} F \mathbf{j}_{\mathrm{Li}^+}$. Substituting \eqref{eq:nernst_planck_flux} and using the local electroneutrality assumption (justified for typical SSB SEs where the ratio of mobile-charge concentration to host-lattice concentration is small) gives the familiar form
\begin{equation}
    \mathbf{i} = -\sigma_{\mathrm{ion}}(c_{\mathrm{Li}^+}) \nabla \phi - z_{\mathrm{Li}^+} F D_{\mathrm{Li}^+} \nabla c_{\mathrm{Li}^+}
    \label{eq:ionic_current_density}
\end{equation}
with effective ionic conductivity $\sigma_{\mathrm{ion}}(c_{\mathrm{Li}^+}) = z_{\mathrm{Li}^+}^2 F^2 D_{\mathrm{Li}^+} c_{\mathrm{Li}^+}/(R\vartheta)$.

Conservation of $\mathrm{Li}^+$ in the SE bulk reads
\begin{equation}
    \frac{\partial c_{\mathrm{Li}^+}}{\partial t} + \nabla \cdot \mathbf{j}_{\mathrm{Li}^+} = 0
    \label{eq:Li_conservation_SE}
\end{equation}
which, combined with charge conservation $\nabla\cdot \mathbf{i} = 0$, yields the standard coupled set
\begin{equation}
\begin{aligned}
    \frac{\partial c_{\mathrm{Li}^+}}{\partial t} & = \nabla \cdot\!\left[D_{\mathrm{Li}^+} \nabla c_{\mathrm{Li}^+} + \tfrac{z_{\mathrm{Li}^+} F}{R\vartheta}\, D_{\mathrm{Li}^+} c_{\mathrm{Li}^+}\, \nabla \phi\right], \\
    0 & = \nabla \cdot \!\left[\sigma_{\mathrm{ion}}(c_{\mathrm{Li}^+}) \nabla \phi + z_{\mathrm{Li}^+} F D_{\mathrm{Li}^+} \nabla c_{\mathrm{Li}^+}\right].
\end{aligned}
    \label{eq:np_coupled_pde}
\end{equation}

Coupling to the interphase layer of \secref{subsec:li_diffusion} occurs at the SE/Li-metal interface through Butler-Volmer kinetics: the interfacial current density imposed by an externally-controlled cell potential drop drives the local $\mathrm{Li}^+$ flux normal to the interface, which is consumed (reduced to neutral Li) on the metal side and supplied as a Dirichlet- or flux-type boundary condition for $c_{\mathrm{Li}^+}$ on the SE side. Continuity of the electrochemical potential $\tilde{\mu}_{\mathrm{Li}} = \mu_{\mathrm{Li}^+} + z_{\mathrm{Li}^+} F \phi$ across the interface closes the coupled system.

In our implementation \eqref{eq:np_coupled_pde} is discretized in MOOSE using a coupled scheme in which both $c_{\mathrm{Li}^+}$ and $\phi$ are nodal variables solved monolithically, with the migration term \eqref{eq:nernst_planck_flux} introduced via the dedicated \texttt{ADNernstPlanckConvection} kernel and the conductivity term via \texttt{ADChargedTransportAnisotropy} for anisotropic SEs.

\subsection{Contact problems during metal growth}
Consider the problem of computing temperature/electric potential distribution
in two regions $\Omega^{(1)}$ and $\Omega^{(2)}$ with conductivity's
$k^{(1)}$ and $k^{(2)}$, which are in contact or separated by a
small gap. The electric potential/temperature $\phi^{m}$ in each
region satisfies 

\begin{equation}
-\nabla\cdot k^{(m)}\nabla\phi^{(m)}=f^{(m)}\in\Omega^{(m)}\label{eq:conduction_eq}
\end{equation}

for $m=1,2$, where $f^{(m)}$ is a specified heat source in body
$m$. The shared regions between the regions is denoted $\Gamma_{c}^{(m)}$
and each region can have Dirichlet conditions on $\Gamma_{D}^{(m)}$
and Neumann conditions on $\Gamma_{N}^{(m)}$ respectively. 

Denoting the flux through the surfaces $\Gamma_{c}^{(m)}$ as $q_{c}^{(m)}$.
If there is not energy stored in the contact region, then the fluxes
through $\Gamma_{c}^{(m)}$ is for example

\begin{equation}
q_{c}^{(1)}=k_{gap}(\phi^{(1)}-\phi^{(2)})\label{eq:interface_1}
\end{equation}

\begin{equation}
q_{c}^{(2)}=k_{gap}(\phi^{(2)}-\phi^{(1)})\label{eq:interface_2}
\end{equation}

or any function of the field variables on either side of the interface.
In the case of linear Butler-Volmer Kinetics, $\phi$ is the electric
potential 

\[
q_{c}^{(1)}=\frac{K}{R_{ct}}(\phi^{(1)}-\phi^{(2)}-\phi_{0}^{(2)}(\sigma))
\]

The weak formulation for each can be written in terms of trial and
test spaces, 

\begin{equation}
\int_{\Omega^{(m)}}(k^{(m)}\nabla\phi\cdot\nabla v^{(m)}-f^{(m)}v^{(m)})dx+\int_{\Gamma_{N}^{(m)}}q_{N}^{(m)}v^{(m)}dS+\int_{\Gamma_{c}^{(m)}}q_{c}^{(m)}v^{(m)}dS=0\label{eq:weakform}
\end{equation}

In a simple situation where the fluxes on both sides are the same, this is written as

\begin{equation}
q_{c}^{(1)}=-q_{c}^{(2)}\label{eq:equalityoffluxes}
\end{equation}

However, in the Battery problem, one side of the interface generally has an electronic flux and other side has an ionic flux, the \eqref{eq:equalityoffluxes} becomes 
\begin{equation}
q_{c}^{(1)}=q_{c}^{(2)}\label{eq:equalityofionicfluxes}
\end{equation}

the weak form for the full problem can be written as 

\begin{equation}
\sum_{m=1}^{2}\left[\int_{\Omega^{(m)}}(k^{(m)}\nabla\phi\cdot\nabla v^{(m)}-f^{(m)}v^{(m)})dx+\int_{\Gamma_{N}^{(m)}}q_{N}^{(m)}v^{(m)}dS\right]+\int_{\Gamma_{c}^{(1)}}q_{c}^{(1)}v^{(1)}dS-\int_{\Gamma_{c}^{(1)}}q_{c}^{(1)}v^{(2)}dS\label{eq:full_weak_form}
\end{equation}

In a typical contact algorithm, we can project the test function
$v^{(2)}$ onto $\Gamma_{c}^{(1)}$ using a smoothing operator $P$
such that the weak form becomes

\begin{equation}
\sum_{m=1}^{2}\left[\int_{\Omega^{(m)}}(k^{(m)}\nabla\phi\cdot\nabla v^{(m)}-f^{(m)}v^{(m)})dx+\int_{\Gamma_{N}^{(m)}}q_{N}^{(m)}v^{(m)}dS\right]+\int_{\Gamma_{c}^{(1)}}q_{c}^{(1)}(v^{(1)}-Pv^{(2)})dS\label{eq:fullwwakform2}
\end{equation}

For non-matching meshes we can introduce a lagrange multiplier $\lambda$
such that $\lambda=q_{c}^{(1)}$, using which the weak form can now
be written as 

\begin{equation}
\sum_{m=1}^{2}\left[\int_{\Omega^{(m)}}(k^{(m)}\nabla\phi\cdot\nabla v^{(m)}-f^{(m)}v^{(m)})dx+\int_{\Gamma_{N}^{(m)}}q_{N}^{(m)}v^{(m)}dS\right]+\int_{\Gamma_{c}^{(1)}}\lambda(v^{(1)}-Pv^{(2)})dS\label{eq:full_weak_lagrange1}
\end{equation}

\begin{equation}
\int_{\Gamma_{c}^{(1)}}\left[\lambda-k_{gap}(\phi^{(1)}-\phi^{(2)})\right]\mu dS\label{eq:full_weaklagrang2}
\end{equation}

where $\mu$ is the test function for the lagrange multiplier. This
provides a general method of creating a contact layer without modeling
the layer itself. $k_{gap}$ can be a function of mechanical constraints
(gap distance, pressure) etc. It can also be a function of spacial
location. 

Example of this method from a MOOSE implementation. ABAQUS implementation
is not certain at this point. We implement \eqref{full_weaklagrang2}
using a mortar FEM method in moose. 

\subsection{Free Energies and chemical potential}
\label{subsec:free_energies}

We start by assuming the free energy of the system is governed by
\begin{equation}
    \hat{\psi}_R(\mathcal{I}_{\mathbf{E}^e}, c_R) = \psi_R^c + \psi_R^e(\mathcal{I}_{\mathbf{E}^e}, c_R)\label{eq:linear_free_energy_decomp}
\end{equation}

\begin{equation}
    \psi_R^c = \mu_0 C_R + R\vartheta c_R \left(\log\left(\frac{\gamma c_R}{c_{R0}}\right) -1\right) \label{eq:chemical_free_energy}
\end{equation}
Here $\psi_R^c$ is the change in chemical free energy due to mixing of Li with the host material. A simple approximation for the chemical free energy is some initial chemical potential of Li, $R$ is the universal gas constant, $\vartheta$ is the temperature, $c_R$ is the concentration of Li in the reference state of the material and $c_{R0}$ is a reference concentration. 

If we assume that the free energy due to elastic deformation of the host material is isotropic, the second term in \eqref{eq:linear_free_energy_decomp} can now be written as 
\begin{equation}
    \psi_R^e(\mathcal{I}_{\mathbf{E}^e}, c_R) = \psi_R^e(\mathbf{E}^e, c_R) = J^g\left[\frac{1}{2}\mathbf{E}^e:\mathcal{C}\mathbf{E}^e\right] \label{eq:isotropic_elastic_free_energy}
\end{equation}
where $\mathcal{C} = 2G\mathcal{I} + (K-\frac{2}{3}G)\mathbf{1}\otimes\mathbf{1}$ is the isotropic elasticity tensor, $K$ is the bulk modulus, $G$ is the shear modulus an $\mathcal{I}$ and $\mathbf{1}$ are the fourth- and second-order identity tensors. The multiplication of $J^g$ in \eqref{eq:isotropic_elastic_free_energy} gives the free energy in the reference space. 

We can now write the total chemical potential of Li as 
\begin{equation}
\mu = \mu^{chem} + \mu^{e} + \mu^{g} \label{eq:chemical_potential_decomp}
\end{equation}
which can now be expanded to 
\begin{equation}
    \mu = \mu_0 + R\vartheta\left(\log\left(\frac{\gamma c_R}{c_{R0}}\right)\right) + \Omega\left[\frac{1}{2}\mathbf{E}^e:\mathcal{C}\mathbf{E}^e\right] - \Omega \mathbf{M}^g:\mathbf{S}^g
\end{equation}
The derivation of the last term is done elsewhere and is not shown here.

\subsection{Numerical Implementation}
\label{subsec:numerical_implementation}

The framework is implemented as a custom application within the MOOSE finite-element platform~\cite{TBD:MOOSE}, which provides parallel mesh management, the Newton-Krylov nonlinear solver, automatic differentiation through the MetaPhysicL library, and adaptive mesh refinement infrastructure. Each physics ingredient introduced in \secref{sec:methodology} is encapsulated as a MOOSE \texttt{Kernel}, \texttt{Material}, \texttt{InterfaceKernel}, or \texttt{BoundaryCondition} object, exposing material parameters and coupling targets through input-file syntax.

\textbf{Spatial discretization.} Plane-strain or axisymmetric ($RZ$) two-dimensional meshes are used throughout the benchmark suite, with linear or quadratic Lagrange elements (\texttt{QUAD4} / \texttt{QUAD9}) for displacement and concentration fields. Mesh files are generated in Gmsh~\cite{TBD:Gmsh} and read via \texttt{FileMeshGenerator}. For problems with sharp interphase boundaries the mesh refinement is concentrated within $\pm 5\,\mu\mathrm{m}$ of the SE/anode interface; uniform refinement (\texttt{uniform\_refine}) is applied for sensitivity studies.

\textbf{Temporal discretization.} An implicit backward-Euler scheme is used for all transient terms, with adaptive time-stepping driven by Newton-iteration-count feedback (\texttt{IterationAdaptiveDT}): on convergence within an optimal-iteration window the next time step is grown by a factor up to $1.5$; on convergence failure the step is cut by $0.65$. A material-time-step postprocessor (\texttt{MaterialTimeStepPostprocessor}) further bounds $\Delta t$ from above by the local creep / diffusion time scale, ensuring stable integration of the viscoplastic and growth response.

\textbf{Coupled solve and scaling.} The momentum, mass-conservation, and charge-conservation residuals are assembled monolithically in a single Newton-Krylov solve. Residual scaling (\texttt{automatic\_scaling = true}) is recomputed each step (\texttt{compute\_scaling\_once = false}) to handle the orders-of-magnitude span between mechanical residuals (in MPa~mm) and electrochemical residuals (in mA~mm). Variable groups (\texttt{scaling\_group\_variables = `disp\_x disp\_y conc'}) keep displacements and concentration on a common scale to avoid spurious dominance of one block.

\textbf{Linear solver.} A direct LU factorization (PETSc \texttt{-pc\_type lu}) with \texttt{superlu\_dist} as the parallel backend is the default. For larger meshes a multigrid preconditioner (\texttt{-pc\_type hypre} with \texttt{boomeramg}) becomes practical; the default direct solver was retained for the benchmark cases reported here for reproducibility.

\textbf{Mortar contact.} The contact formulation of \secref{subsec:contact} is implemented via MOOSE's mortar method, in which an auxiliary Lagrange-multiplier variable $\lambda$ is defined on the secondary side of the interface (\texttt{frictionless\_secondary\_subdomain}). Contact conditions $\lambda - k_{\mathrm{gap}}(\phi^{(1)} - \phi^{(2)}) = 0$ are imposed weakly via the \texttt{PressureLMConductanceConstraint} object. This avoids needing matched meshes across the SE/anode interface and supports both normal and tangential constraints.

\textbf{Reproducibility.} All inputs reported in \secref{sec:benchmarks} and \secref{sec:plating_stripping} are version-controlled within the \texttt{ecm\_test} application repository and reproducible from a fresh clone after running the benchmarks via the supplied scripts. Adaptive parameters were not hand-tuned per case; the same default time-stepper / solver settings (with the exception of mesh resolution) were used across the suite.

\section{Benchmark Tests}
\label{sec:benchmarks}

The developed framework was tested step by step by following the procedure: Plastic/creep behavior of Li metal material $\to$ Uniform Li metal growth in the interphase layer $\to$ Li metal growth incorporating Li diffusion in the interphase layer $\to$ Metal growth between conformal contact of Li-metal anode and the SE $\to$ Incorporating Li-ion migration in the SE.

\subsection{Plastic/creep behavior of Li metal material}
\label{subsec:plastic_creep}

The viscoplastic-creep response of Li metal is the prerequisite for any quantitative model of plating-driven stress, since at room temperature Li approaches half its melting point and its flow stress at relevant strain rates is comparable to typical SE fracture toughness. We adopt the Anand-type rate-dependent model with a single internal hardening variable $s$ (the Anand resistance), governed by a hardening-saturation evolution law. The active rate equations and parameter values are inherited from the Li-metal extension proposed by~\cite{TBD:LiAnand}; the implementation is contained in the \texttt{ADIsoTropicHyperVisco} family of materials in \texttt{ecm\_test/src/materials/}.

The benchmark exercises the constitutive law in two protocols: (i) a constant true-strain-rate sweep at room temperature (\figref{fig:differentStrainRate}) and (ii) a temperature sweep at a single strain rate (\figref{fig:differentTemperature}). Both protocols are run as one-element axisymmetric ($RZ$) tension/compression tests via \texttt{benchmarks/benchmark1\_metalCreep/tension\_constTrueStrainRate.i} and parametric variants. The simulated curves bracket the available experimental data (cited in~\cite{TBD:LeSar,TBD:LePage,TBD:Masias}) across both ranges and reproduce the rate-sensitivity exponent and activation-energy controlled temperature dependence within experimental scatter.

\begin{figure}[h!]
    \centering
    \includegraphics[width=9cm]{StrainRate.png}
    \caption{True stress-strain response of Li metal at $T = 298\,\mathrm{K}$ for steady-state true-strain rates between $4\times 10^{-5}\,\mathrm{s}^{-1}$ and $2\times 10^{-2}\,\mathrm{s}^{-1}$. The rate sensitivity reproduces the Anand-model predictions and matches the experimentally-measured rate-exponent in this strain-rate window. Simulation source: \texttt{benchmark1\_metalCreep/tension\_constTrueStrainRate.i}.}
    \label{fig:differentStrainRate}
\end{figure}

\begin{figure}[htp]
    \centering
    \includegraphics[width=10cm]{temperature.png}
    \caption{True stress-strain response at fixed strain rate $\dot{\varepsilon} = 4\times 10^{-5}\,\mathrm{s}^{-1}$ for several temperatures spanning typical battery operating conditions. The activation-energy controlled softening with temperature follows Arrhenius scaling and matches the experimentally-reported temperature trend. Simulation source: \texttt{benchmark1\_metalCreep/tension\_constTrueStrainRate.i} (with parametric temperature).}
    \label{fig:differentTemperature}
\end{figure}

The agreement across both rate and temperature axes establishes the constitutive parameter set used for Li in all subsequent benchmarks (\subsecref{Free-growth} through \secref{sec:plating_stripping}). Where solid-state Li-metal cell measurements deviate from these isotropic constitutive predictions, the deviation is traceable to grain-boundary or interfacial effects not captured at the continuum level here; quantification of those effects is left for future work.



\subsection{Free growth of Li metal in the interphase layer}
\label{subsec:Free-growth}

Uniform Li metal growth was first studied to test the convergence behavior, the volumetric growth response, the effect of initial interlayer thickness, and the response on curved deposition surfaces. The reference project is \texttt{benchmarks/benchmark2\_freeGrowth/}.

\subsubsection*{Uniform normal growth on flat surface}
The simplest verification case prescribes a spatially-uniform concentration field in a flat-geometry interlayer block; the growth tensor of \eqref{eq:Jg_constitutive} should produce uniform vertical extension. Reference input: \texttt{01\_flat\_uniformSwell.i}.

\begin{figure}[h!]
    \centering
    \framebox[12cm]{\rule[-2.5cm]{0pt}{5cm}\textit{TBD: deformed mesh + vertical-displacement profile for flat uniform growth}}
    \caption{Flat-geometry interlayer under spatially-uniform concentration ramp. Vertical displacement scales linearly with $\Omega(c_R - c_{R0})$ as predicted by \eqref{eq:Jg_constitutive}, confirming the kinematic implementation.}
    \label{fig:freegrowth_flat}
\end{figure}

\subsubsection*{Uniform normal growth on inclined and cosine-shape surface}
The same uniform-concentration ramp is applied on inclined and cosine-shape surfaces. With areal-growth weights $\alpha_1 = 0$, $\alpha_2 = \alpha_3 = 1/2$ the deposition direction tracks the surface normal even as the surface is tilted or curved. Reference inputs: \texttt{inclined.i}, \texttt{cosine.i}.

\begin{figure}[h!]
    \centering
    \framebox[12cm]{\rule[-2.5cm]{0pt}{5cm}\textit{TBD: deformation snapshots for inclined and cosine surfaces, color = displacement magnitude}}
    \caption{Areal growth on (left) an inclined surface and (right) a cosine-shape surface. The deposited Li thickness is uniform along the normal direction, demonstrating the principal-vector transformation of \eqref{eq:principal_vector_transformation} under non-trivial reference geometry.}
    \label{fig:freegrowth_curved}
\end{figure}

\subsubsection*{Effect of initial thickness of the interlayer}
A thickness sweep ($h_0 \in \{1, 2, 5, 10\}\,\mu\mathrm{m}$) reveals the role of the interlayer as a stress buffer: thinner interlayers transfer more of the growth strain into the bonded SE/anode constraint, raising the contact-pressure response. Reference input: \texttt{swell.i} (with parametric thickness).

\begin{figure}[h!]
    \centering
    \framebox[12cm]{\rule[-2.5cm]{0pt}{5cm}\textit{TBD: contact pressure vs.\ deposited Li volume for four interlayer thicknesses}}
    \caption{Contact pressure at the SE/anode interface as a function of deposited Li volume for four initial interlayer thicknesses $h_0$. Thinner interlayers exhibit steeper pressure build-up, consistent with reduced compliance.}
    \label{fig:freegrowth_thickness}
\end{figure}

When $\mathrm{Li}^+$ deposits between the SE and Li-metal interface, we assume the newly deposited Li atoms undergo diffusion within the interlayer. Notably, this Li diffusion is along the principal direction, following the growth direction.

\subsubsection*{Li diffusion-induced growth on inclined and cosine-shape surface}
Replacing the prescribed concentration ramp with a flux boundary condition activates the diffusion-coupled growth: Li injected at the SE-side boundary diffuses into the interlayer through the anisotropic mobility tensor of \eqref{eq:reference_Li_flux}, and the local concentration drives local growth via $\dot{\epsilon}^g = \Theta(c_R)\dot{c}_R$. Reference inputs: \texttt{02\_flat\_diffusionSwell.i}, \texttt{diffusion.i}.

\begin{figure}[h!]
    \centering
    \framebox[12cm]{\rule[-2.5cm]{0pt}{5cm}\textit{TBD: snapshots of concentration and displacement on cosine surface during diffusion-driven growth}}
    \caption{Diffusion-driven growth on a cosine-shape surface. (a) Li concentration $c_R$ at three times shows the diffusion front advancing into the interlayer. (b) Surface displacement co-located with the concentration front confirms the constitutive coupling between $c_R$ and $J^g$.}
    \label{fig:freegrowth_diffusion}
\end{figure}

\subsubsection*{Effect of Li diffusivity in the interlayer}
A diffusivity sweep ($D \in \{1, 10, 100, 1000\}\,\mathrm{nm}^2/\mathrm{s}$) at fixed flux quantifies the regime transition between diffusion-limited (slow $D$, sharp concentration front, slow surface advance) and reaction-limited (fast $D$, near-uniform concentration, surface advance set by boundary flux) growth. Reference input: \texttt{02\_flat\_diffusionSwell\_Ash.i} swept over diffusivity.

\begin{figure}[h!]
    \centering
    \framebox[12cm]{\rule[-2.5cm]{0pt}{5cm}\textit{TBD: surface advance rate vs.\ diffusivity, log-log axes}}
    \caption{Surface advance rate as a function of interlayer Li diffusivity at fixed boundary flux. The two asymptotes (diffusion-limited at low $D$, reaction-limited at high $D$) bracket a transition regime in which both effects matter.}
    \label{fig:freegrowth_diffusivity}
\end{figure}

\subsection{Confined growth of Li metal between anode and SE}
\label{subsec:confined_growth}

The confined-growth benchmark embeds the interlayer between a rigid SE and a Li-metal anode in conformal contact, so growth at the SE/interlayer interface is resisted by mechanical compression on the anode side. The contact is mediated by the mortar-Lagrange-multiplier formulation of \secref{subsec:contact}. Five reference inputs span (i) mechanics-only baseline, (ii) coupled growth, and (iii–v) three surface geometries (circle, cosine, incline). Project: \texttt{benchmarks/benchmark3\_confinedGrowth/}.

\begin{figure}[h!]
    \centering
    \framebox[14cm]{\rule[-3cm]{0pt}{6cm}\textit{TBD: contact pressure and deformation for three surface geometries}}
    \caption{Confined growth in three SE-surface geometries: (a) circular dimple (\texttt{test\_circle.i}), (b) cosine corrugation (\texttt{test\_cosine.i}), (c) inclined ramp (\texttt{test\_incline.i}). The contact pressure distribution shows local stress concentration at curvature extrema, predicting the locations most susceptible to SE fracture or Li penetration. The mechanics-only baseline (\texttt{01\_conformalContact\_mech.i}) and the coupled growth case (\texttt{02\_conformalContact\_swell.i}) bracket the role of growth in contact-pressure evolution.}
    \label{fig:confined_growth}
\end{figure}

The benchmark establishes that growth in conformal contact builds compressive pressure that scales with the imposed Li flux and the anode-side compliance, and that surface curvature focuses pressure into hot-spots whose magnitudes are predicted quantitatively by the framework.

\subsection{Full coupling with Li-ion migration in the SE}
\label{subsec:full_coupling}

The fully-coupled benchmark combines the diffusion-driven growth in the interphase (\subsecref{Free-growth}) with $\mathrm{Li}^+$ Nernst-Planck migration in the SE (\secref{subsec:li_migration_se}) and Butler-Volmer interfacial kinetics. Galvanostatic boundary conditions on the SE outer surface drive a $\mathrm{Li}^+$ flux that, after charge transfer at the SE/interphase interface, supplies the neutral Li concentration that in turn drives growth. The reference inputs are \texttt{benchmarks/benchmark3\_growthDiffusion/benchmark3\_growthDiffusion.i} (single SE/interphase pair) and \texttt{benchmarks/benchmark4\_fullCoupled/full\_model.i} (full SE/interphase/Li-metal stack with mortar contact).

\begin{figure}[h!]
    \centering
    \framebox[14cm]{\rule[-3cm]{0pt}{6cm}\textit{TBD: $\mathrm{Li}^+$ concentration profile in SE and Li concentration in interphase, snapshots at three times}}
    \caption{Steady-state coupling between $\mathrm{Li}^+$ flux in the SE and Li metal growth in the interphase under galvanostatic plating. (a) $\mathrm{Li}^+$ concentration profile in the SE shows the depletion / accumulation gradient set by the imposed current. (b) Li concentration in the interphase shows the corresponding accumulation that drives growth via the constitutive law of \eqref{eq:Jg_constitutive}. (c) Resulting interface displacement vs.\ time tracks the volumetric growth rate $\dot{\epsilon}^g$.}
    \label{fig:full_coupling}
\end{figure}

The benchmark validates that the coupled system reaches the analytical steady state in which the $\mathrm{Li}^+$ migration flux equals the rate of Li-metal growth at the interface, modulo the kinetic barrier set by the Butler-Volmer exchange-current density. Beyond steady-state validation, the case provides the geometric configuration used in \secref{sec:plating_stripping} for plating / stripping studies, and demonstrates that the framework supports the full coupling chain at industrially-relevant cell geometries.

\section{Li Metal Plating/Stripping under Different Conditions}
\label{sec:plating_stripping}

The benchmarks of \secref{sec:benchmarks} establish that the framework reproduces canonical viscoplastic, growth, contact, and migration responses individually. We now combine them to study Li metal plating and stripping in representative solid-state cell architectures, where the coupling between Li-ion transport in the SE, Li metal growth at the interphase, and contact pressure at the SE/anode interface dictates whether the cell cycles stably or fails by dendrite penetration or void formation. Five sub-cases cover the parameter regimes most relevant to current experimental observations.

\subsection{Galvanostatic plating with a bonded interlayer (base case)}
\label{subsec:plating_base}

We first consider a \emph{bonded} configuration in which a thin Li reservoir, a porous Li interlayer, and a ceramic SE are in tied contact at all interfaces. Galvanostatic plating is imposed by applying a constant Li-ion flux at the SE outer boundary; the resulting Li accumulation at the SE/interlayer interface drives growth through the kinematic framework of \secref{sec:methodology}. The reference input is \texttt{benchmarks/plating\_and\_stripping/interlayer\_left\_bond\_pore\_left\_bond.i}.

\begin{figure}[h!]
    \centering
    % \includegraphics[width=14cm]{TBD_plating_base.png}
    \framebox[14cm]{\rule[-3cm]{0pt}{6cm}\textit{TBD: contact-pressure and Li-concentration evolution during galvanostatic plating}}
    \caption{Evolution of (a) contact pressure at the SE/Li interface and (b) Li concentration in the porous interlayer during galvanostatic plating at $J = J_{\mathrm{plating}}$ for the bonded-interlayer base case. The bonded constraint forces the interlayer to grow vertically, generating compressive contact pressure that builds monotonically until plastic relaxation in the reservoir Li bounds the steady-state stress.}
    \label{fig:plating_base}
\end{figure}

The base case demonstrates the framework's ability to track plating-induced stress build-up self-consistently, without prescribing the contact pressure as an external input. As a sanity check we recover the analytical asymptote $\sigma_\infty \approx \mu_{\mathrm{Li}} J_{\mathrm{plating}} t / h_{\mathrm{interlayer}}$ in the limit of small elastic strain and no plastic dissipation; deviations from this limit at later times quantify the role of viscoplastic flow in stress relaxation.

\subsection{Plating into a porous interlayer with stress-driven equilibration}
\label{subsec:plating_porous}

The second case relaxes the bonded constraint on the interlayer side and allows porosity-mediated equilibration: Li deposited at the SE/interlayer interface can redistribute into the porous structure via stress-driven creep before contributing to interface contact pressure. We compare two variants: \texttt{interlayer\_left\_bond\_pore\_left\_bond\_w\_equilibriation.i} (with equilibration) and the base case (without). The reference \texttt{interlayer\_left\_bond\_pore\_left\_bond\_mech.i} provides the mechanics-only counterpart that omits Li-ion migration, isolating the role of mechanical equilibration alone.

\begin{figure}[h!]
    \centering
    \framebox[14cm]{\rule[-3cm]{0pt}{6cm}\textit{TBD: comparison of contact pressure and densification with vs.\ without equilibration}}
    \caption{Contact pressure at the SE/Li interface (left) and porosity profile through the interlayer (right) for the same plating current, with and without stress-driven Li equilibration in the porous layer. Equilibration redistributes plated Li away from the interface, lowering peak contact pressure but densifying the interior of the interlayer.}
    \label{fig:plating_porous}
\end{figure}

The comparison quantifies the trade-off between interface stress and bulk densification: a fast equilibration time scale lowers the interfacial pressure (favorable for avoiding SE fracture) at the cost of interior densification (which can shut off transport in continued cycling). This trade-off is one of the central physical questions a unified ECM framework can address that fluid-only or mechanics-only models cannot.

\subsection{Stripping under tension}
\label{subsec:stripping}

Reversing the current density runs the system in stripping mode. Two outcomes are physically possible: a uniform retreat of the Li/interlayer interface (the desired mode) or void nucleation at the SE/Li interface when the local tensile stress exceeds the metal's plastic-flow threshold and Li cannot be supplied fast enough. The reference input is \texttt{benchmarks/plating\_and\_stripping/interlayer\_left\_bond\_pore\_left\_bond\_strip.i}.

\begin{figure}[h!]
    \centering
    \framebox[14cm]{\rule[-3cm]{0pt}{6cm}\textit{TBD: stripping-induced void formation at the SE/Li interface}}
    \caption{Snapshots during galvanostatic stripping showing (a) interface displacement from the reference SE position and (b) Li chemical-potential field in the interlayer. At sufficiently high stripping current density the interface develops localized retreat (void precursor) where the local Li flux from the interlayer cannot sustain the imposed boundary current.}
    \label{fig:stripping}
\end{figure}

The simulation reproduces the experimentally-observed onset of voiding under high-current stripping and provides a quantitative criterion in terms of the ratio between imposed current density and the maximum sustainable Li flux from the porous interlayer.

\subsection{Validation against McMeeking-type experiments}
\label{subsec:mcmeeking}

To anchor the framework against an external benchmark we reproduce the experimental protocol from McMeeking and co-workers' studies of Li-metal SSB cells with cathode-side stack pressure. The setup consists of a ceramic SE between a Li-metal anode and a composite cathode, with a controlled stack pressure imposed at the cathode-side current collector. Reference inputs span a sweep of stack-pressure conditions: \texttt{benchmarks/mcmeeking\_paper\_results/test\_w\_cathode.i} (cathode loaded), \texttt{test\_w\_Ag.i} (Ag interlayer at the anode side), \texttt{test\_w\_equil\_metal.i} (equilibrated metal anode), and \texttt{interface\_grad.i} (with explicit interface gradients).

\begin{figure}[h!]
    \centering
    \framebox[14cm]{\rule[-3cm]{0pt}{6cm}\textit{TBD: comparison vs.\ McMeeking experimental data — interface pressure and current density evolution}}
    \caption{Predicted (lines) vs.\ experimentally-reported (markers) interface contact pressure as a function of cycling time for three stack-pressure conditions, reproducing the protocol of McMeeking et al. The framework captures both the early transient (elastic build-up) and the late-time creep relaxation observed experimentally.}
    \label{fig:mcmeeking}
\end{figure}

Quantitative agreement within experimental uncertainty validates the constitutive parameters chosen for the Li-metal viscoplastic response (\secref{subsec:plastic_creep}) when applied to a more complex cell-level test, and provides a foundation for using the framework predictively for cell architectures not yet explored experimentally.

\subsection{Dendrite onset and propagation via phase-field coupling}
\label{subsec:dendrite}

Finally, we demonstrate the framework's ability to capture dendrite formation by coupling the ECM kernels to a phase-field representation of the Li/SE interface. The phase-field variable $\phi$ tracks the position of the Li metal within the interphase, with its evolution driven by the local electrochemical overpotential modulated by mechanical stress. The reference inputs are \texttt{benchmarks/dendrite\_growth/DendriteGrowth.i} (the canonical case) and \texttt{phaseField\_coupling\_pot.i} (with explicit potential coupling).

\begin{figure}[h!]
    \centering
    \framebox[14cm]{\rule[-3cm]{0pt}{6cm}\textit{TBD: dendrite morphology evolution under galvanostatic plating}}
    \caption{Phase-field snapshots showing dendrite initiation and propagation under galvanostatic plating with surface-energy anisotropy. Left: morphology at $t = t_1$ (initial perturbation amplification). Center: $t = t_2$ (filament extension). Right: $t = t_3$ (impingement on SE counter-electrode). Color denotes the phase-field variable $\phi$ ($\phi=1$ is Li metal, $\phi=0$ is SE).}
    \label{fig:dendrite}
\end{figure}

Coupling the phase-field with the mechanical stress field (rather than running phase-field alone) is essential for capturing the experimentally-observed sensitivity of dendrite morphology to stack pressure: at higher stack pressure compressive stress at the dendrite tip suppresses Li deposition there, blunting the filament; at low stack pressure stress relief at the tip accelerates extension. The framework reproduces this competition self-consistently through the Mandel-stress coupling of \secref{subsec:free_energies}.


\section{Conclusions}
\label{sec:conclusions}

We have presented a unified electro-chemo-mechanical finite-element framework for solid-state lithium-metal batteries. The framework rests on three coupled physical pillars: (i) a multiplicative kinematic decomposition of the deformation gradient $\mathbf{F} = \mathbf{F}^e\mathbf{F}^p\mathbf{F}^g$ that consistently treats elastic, viscoplastic, and growth-driven deformation; (ii) anisotropic Li diffusion in the interphase coupled to Nernst-Planck migration of Li${}^+$ in the solid electrolyte through interfacial Butler-Volmer kinetics; and (iii) mortar-method contact between the SE, interlayer, and Li metal anode that supports both bonded and frictionless interfaces with thermal/electrical conduction analogues. A series of benchmarks --- viscoplastic creep, free swelling, confined growth, fully-coupled migration, and Li plating/stripping under multiple cell architectures --- demonstrate that the framework captures both individual physics responses and their coupled behavior. Validation against McMeeking-type cell-level experiments and reproduction of dendrite-onset phenomenology suggest that the framework is sufficiently general to inform the design of next-generation SSB architectures.

Future work will extend the framework to (a) heterogeneous SE microstructures with grain-boundary transport, (b) thermal coupling for joule-heating effects under fast charging, and (c) statistical analysis of dendrite-onset thresholds via parametric sweeps over interfacial roughness and stack-pressure conditions.

\end{document}
